\chapter{Levels, Objects, and Consequence}
\label{ch:levels}

\section{Formal logic and a logical system}
\label{sec:levels-01}

\begin{definition}[Formal language]
A \emph{formal language} is a set of finite expressions determined from a
specified alphabet by a specified grammar.  A formal language is an abstract
syntactic object.  The grammar determines which finite strings count as
expressions of each grammatical category.
\end{definition}

\begin{definition}[Calculus]
A \emph{calculus} is a specified collection of initial derivations and
inference rules that inductively determines a derivability relation.  The
standing of a calculus is syntactic: it concerns expressions and rules, not
their interpretations.
\end{definition}

\begin{definition}[Semantics]
A \emph{semantics} for a formal language specifies the complete semantic data
permitted for that language and a satisfaction relation between those data and
formulas.  The data may be a propositional valuation, a first-order structure
with a variable assignment, a modal structure with a world, or another tuple
declared explicitly by the semantics.  The satisfaction relation supplies the
truth conditions of the language.
\end{definition}

\begin{definition}[Logical system]
A \emph{logical system} is an ordered triple
\[
  \tuple{\Lang,\mathcal C,\mathcal S},
\]
where \(\Lang\) is a formal language, \(\mathcal C\) is a calculus for
\(\Lang\), and \(\mathcal S\) is a semantics for \(\Lang\).  The term is used
only when all three components have been identified.
\end{definition}

The ordered triple prevents a common error.  A symbol such as \(\derives\)
does not name one universal relation: it abbreviates derivability in a
particular calculus.  Likewise, \(\entails\) depends on a language and its
specified semantic data.  When more than one calculus or semantics is active,
this document writes subscripts explicitly.

\begin{example}
For classical propositional logic, \(\Lang\) contains propositional formulas,
\(\mathcal C\) may be a natural-deduction calculus, and \(\mathcal S\) uses
two-valued valuations.  Replacing natural deduction by a sequent calculus
changes \(\mathcal C\) but not necessarily \(\Lang\) or \(\mathcal S\).
Replacing two-valued valuations by three-valued matrices changes
\(\mathcal S\), and the resulting semantic consequence relation must be
defined again.
\end{example}

